% ex: ts=2 sw=2 sts=2 et filetype=tex
% SPDX-License-Identifier: CC-BY-SA-4.0

\begin{frame}
    \frametitle{Contenido}
    \tableofcontents
\end{frame}

\section{Introducción}

\begin{frame}[c]{Definición de algoritmo}
  \begin{block}{Definición}
    \begin{itemize}
      \item Conjunto ordenado y finito de operaciones que permite hallar
            la solución de un problema\footnote{Real academia española}.
      \pausa
      \item Conjunto de pasos, acciones o instrucciones necesarios para
            lograr un resultado o resolver un problema.
      \pausa
      \item Secuencia ordenada de pasos exentos de ambigüedad tal que, al
            llevarse a cabo con fidelidad, dará como resultado que se realice
            la tarea para la que se ha diseñado en un tiempo finito.
    \end{itemize}
    \pausa
    \begin{center}
    \begin{tikzpicture}[node distance = 2cm]
      \node (problema) [rectangle, fill=none]{Problema};
      \node (flecha) [single arrow, minimum height=1.5cm, right of=problema, fill=cyan!50]{};
      \node (solucion) [rectangle, fill=none, right of=flecha] {Solución};
    \end{tikzpicture}
    \end{center}
  \end{block}
\end{frame}

\begin{frame}[c]{Propiedades de un algoritmo}
  \begin{description}
    \item[E/S] Tiene definidas las entradas que se requieren, así como las
      salidas que se deben producir.
    \pausa
    \item[Preciso] Se debe indicar sin ambigüedades el orden exacto de los
      pasos a seguir y la manera en la que éstos deben realizarse.
    \pausa
    \item[Determinista] para un mismo conjunto de datos proporcionado como
      entrada, el algoritmo debe producir siempre el mismo resultado.
    \pausa
    \item[Finito] un algoritmo siempre termina después de un número finito
      de pasos.
    \pausa
    \item[Efectivo] cada paso o acción se debe poder llevar a cabo en un
      tiempo finito y se debe lograr el efecto que se desea o espera.
  \end{description}
\end{frame}

\begin{frame}[c]{Análisis de algoritmos}
  \begin{itemize}
    \item El análisis de algoritmos pretende descubrir si estos son o no
      eficaces.
    \pausa
    \item Establece una comparación entre los mismos con el fin de saber
      cuál es el más eficiente.
    \pausa
    \item No importa que cada uno de los algoritmos de estudio sirva para
      resolver el mismo problema.
  \end{itemize}
\end{frame}

\begin{frame}[c]{Eficiencia de los algoritmos}

  Para saber qué tan eficiente es un algoritmo hacemos las preguntas:
  \vspace{\baselineskip}

  \pausa
  \begin{itemize}
    \item ¿Cuánto tiempo ocupa?
    \pausa
    \item ¿Cuánta memoria ocupa?
  \end{itemize}
  \pausa
  \begin{block}{Nota:}
    El tiempo que requiere un algoritmo para resolver un problema está
    en función del tamaño $n$ del conjunto de datos para procesar: $T(n)$
  \end{block}
\end{frame}

\begin{frame}[c]{Tipos de estudio del tiempo}
  \begin{description}
    \item[A priori] se obtiene una medida teórica, es decir, una función
      que acota (por arriba o por abajo) el tiempo de ejecución del
      algoritmo para unos valores de entrada dados.
    \pausa
    \item[A posteriori] se obtiene una medida real, es decir, el tiempo
      de ejecución del algoritmo para unos valores de entrada dados en
      determinada computadora.
  \end{description}
  \pausa
  \begin{alertblock}{Nota:}
    El tiempo de ejecución $T(n)$ de un algoritmo para una entrada de
    tamaño $n$ no debe medirse en segundos, milisegundos, etc.
  \end{alertblock}
\end{frame}


\section{La eficiencia de un algoritmo}

\begin{frame}[c]{Principio de invarianza}
  \begin{block}{Principio de invarianza}
    Dado un algoritmo y dos implementaciones suyas $I_1$ e $I_2$, que tardan
    $T_1(n)$ y $T_2(n)$ respectivamente, el \textbf{Principio de invarianza}
    afirma que existe una constante real $c > 0$ y un número natural $n_0$
    tales que para todo $n \geq n_0$ se verifica que:
    \[
      T_1(n) \leq cT_2(n)
    \]
  \end{block}
\end{frame}

\begin{frame}[c]{En otras palabras}
  \begin{itemize}
    \item El tiempo de ejecución de un algoritmo es asintóticamente de
      \textit{del orden de } $T(n)$ si existen dos constantes reales
      $c$ y $n_0$ y una implementación $I$ del algoritmo tales que el
      problema se resuelve en menos de $cT(n)$, para toda $n > n_0$.
  \end{itemize}
\end{frame}

\begin{frame}[c]{Influye más el tipo de función}
  \begin{table}[]
  \begin{tabular}{rrcl}
    \textbf{n} &  \textbf{Algoritmo 1} &   &  \textbf{Algoritmo 2} \\
       & $T_1(n) = 10^6n^2$  &   & $T_2(n) = 5n^3$  \\ 
         2 & $10^6 \times 4 = 4'000,000$ & > & $5 \times 8 = 40$ \\
       200 & $10^6 \times 40,000 = 4 \times 10^{10}$ & > & $5 \times 8 \times 10^6 = 4 \times 10^7$  \\
       200,000 & $10^6 \times 4 \times 10^{10} = 4 \times 10^{16}$ & = & $5 \times 8 \times 10^{15} = 4 \times 10^{16}$  \\
       2'000,000 & $10^6 \times 4 \times 10^{12} = 4 \times 10^{18}$ & < & $5 \times 8 \times 10^{18} = 4 \times 10^{19}$  \\
  \end{tabular}
  \end{table}

  \vspace{\baselineskip}
  A partir de cierto valor de $n$, la función cúbica es siempre mayor a
  pesar de que la constante es mucho menor a ésta.
\end{frame}


\section{El tiempo de ejecución $T(n)$ de un algoritmo}

\begin{frame}[c]{El tiempo de ejecución}
  Para medir $T(n)$ usamos el número de operaciones elementales. Una operación
  elemental puede ser:
  \pausa
  \begin{itemize}
    \item Operación aritmética.
    \pausa
    \item Asignación a una variable.
    \pausa
    \item Llamada a una función.
    \pausa
    \item Retorno de una función.
    \pausa
    \item Comparaciones lógicas (con salto).
    \pausa
    \item Acceso a una estructura (arreglo, matriz, lista ligada…).
  \end{itemize}

  \vspace{\baselineskip}
  Se le llama \textbf{tiempo de ejecución}, no al tiempo físico, sino al
  número de operaciones elementales que se llevan a cabo en el algoritmo.

\end{frame}
